% ---------------------------------------------------------------------------
% Supervisor meeting, 2026-11-03: what has been done to speed up the sampler,
% what it gained, what did not help, and the options from here.
%
% Every number is taken from softmax/DEVLOG.txt (entries 2026-10-22 to
% 2026-11-02) or from the files named next to it. Companion to
% one_pass_move.tex (the one-pass move in detail) and zero_polymer_recap.tex.
%
% Build:  cd softmax/docs && pdflatex speedup_status.tex && pdflatex speedup_status.tex
% ---------------------------------------------------------------------------
\documentclass[11pt,a4paper]{article}

\usepackage[T1]{fontenc}
\usepackage[utf8]{inputenc}
\usepackage[margin=2.4cm]{geometry}
\usepackage{amsmath,amssymb}
\usepackage{booktabs}
\usepackage{array}
\newcolumntype{R}[1]{>{\raggedright\arraybackslash}p{#1}}
\usepackage[hidelinks]{hyperref}
\sloppy

\newcommand{\dt}{\Delta t}
\newcommand{\Hd}{H_d}
\newcommand{\file}[1]{\texttt{#1}}

\title{Speeding up the sampler: status and options\\
       \large Notes for the supervisor meeting}
\author{Marco Sanchini \\ \small\texttt{my\_extendedSampling} / \texttt{softmax}}
\date{2026--11--03}

\begin{document}
\maketitle

\begin{abstract}
\noindent
Goal: make the gradient-informed sampler on the zero\_polymer reference
(influenza HA0, $L=566$) fast enough to reach equilibrium at and below the
folding crossover. One change has paid off. The \emph{one-pass move} computes
the exact energy and the gradient for all sites with one forward and one
backward pass through ESM3. It costs 0.153\,s per move instead of 0.365\,s,
a $2.38\times$ speed-up with the same per-move statistics and an exact
Metropolis--Hastings kernel. Three other ideas were measured and give little
or nothing: running two chains on the GPU at once ($-7\%$), mutating several
sites per move ($\le1.16\times$, only at high $T$), and replica exchange (no
gain per move at an affordable length). The limit now is acceptance, not cost:
at $T\le15$ only 0.5--6\% of moves are accepted. The most promising next step
is an exact single-site move that evaluates all 19 alternatives in one batched
forward (Metropolized Gibbs), worth up to $\sim12\times$ at low $T$ if batching
is efficient. It is not built yet, and it removes the gradient from the
proposal, which is the main decision for this meeting.
\end{abstract}

\tableofcontents

% ===========================================================================
\section{Where the time goes}

A move picks a site uniformly, proposes a new residue from a softmax over the
gradient-predicted energy change (drift $-\dt^2\nabla U/2M$, noise
$\dt\sqrt{T/M}$; $\dt/\sqrt{T}$ is fixed across $T$, so the informedness is
the same at every temperature), and accepts with the Metropolis--Hastings
ratio. The energy $U$ is ESM3 cross-entropy. Every cost is GPU time on the
RTX 4500 Ada (24\,GB, ``astaroth''), and the GPU is fully busy during a move.

\begin{center}
\begin{tabular}{lccc}
\toprule
move & network passes per move & s/move & $10^5$ moves \\
\midrule
site\_soft (before)  & 3 forward + 2 backward & 0.365 & 10.1\,h \\
one-pass (now)       & 1 forward + 1 backward & 0.153 & 4.3\,h \\
\bottomrule
\end{tabular}
\end{center}

\noindent
Why this matters: at $T=100$ two reference-started and two random-started
chains needed $\sim$140\,000 one-pass moves ($\sim$6\,h each) to agree on
$\Hd$ ($\hat R_U=1.48$, $\hat R_{\Hd}=1.13$). At $T=32$ they still disagree
after 140\,000 ($\hat R_U=5.25$). At $T=12$ a single reference-started chain
needed $\sim$400\,000 moves ($\sim$17\,h) to reach a quasi-plateau. With
single-site moves at acceptance 0.08, touching each of the 566 sites about
$\ln 566$ times already takes $\sim$45\,000 moves.

% ===========================================================================
\section{What has been done}

\subsection{One-pass move: $2.38\times$ per move (done, in production)}

Two changes (details in \file{one\_pass\_move.tex}):
\begin{enumerate}
\item The gradient is taken with respect to the one-hot probability matrix at
  the exact sequence, so the forward pass that gives the exact $U$ also gives
  the gradient.
\item The gradient is computed for all sites at once (same cost as for one
  site). The gradient at the candidate becomes the gradient of the next move
  if the candidate is accepted, so each move needs one forward + one backward.
\end{enumerate}
The gradient changes scale but not its ranking of residues, so $\dt$ is
rescaled by a computable factor $c=p_{\mathrm{off}}/T_{\mathrm{sftm}}=
4.535\times10^{-4}$. The kernel is exact for any gradient, including the
non-bit-reproducible bf16 backward.

\begin{center}
\small
\begin{tabular}{R{1.1cm}R{5.6cm}R{7.2cm}}
\toprule
test & what & result \\
\midrule
T0 & gradient vs finite differences (fp64) & median rel.\ error $4$--$5\times10^{-6}$ \\
T1 & $U$ from the grad-enabled forward & equal to the exact $U$ bit for bit (50/50) \\
T2 & one-pass vs site\_soft gradient (fp32) & cosine $\ge0.999997$, scale $1.001\,c$ \\
T4 & Boltzmann on a 2-site restricted problem & KL/TV $0.0154/0.0527$ (site\_soft $0.0169/0.0569$), both $\propto$ MC noise; acc.\ 0.305 vs 0.306 \\
T5 & informedness vs $\dt$ & $p_{\mathrm{match}}$ agrees site by site to $\sim$1\% \\
T6 & cost & 0.153 vs 0.365\,s/move; peak memory 6350\,MB in both \\
T7a & transient from the reference, $T=32,100,316$ & acceptance agrees ($|z|\le1.4$); $\Hd$, $U$ within a few \% \\
\bottomrule
\end{tabular}
\end{center}
Toy-energy checks (exact enumeration, $L=4,5$, $N=1$--3): detailed balance to
$3\times10^{-18}$; a deliberately wrong variant was detected (TV 0.050 vs
0.008).

\subsection{Two chains on the GPU at once: no gain}

Two chains together run at 0.329\,s/move each (one-pass) or 0.79\,s/move
(site\_soft): a move every 0.164\,s instead of 0.153\,s, i.e.\ 7--8\% slower
than running them one after the other. The scripts now run one job at a time.

\subsection{Several mutations per move ($n_{\mathrm{sites}}$ scan): at most $1.16\times$}

From the $T=32$ and $T=100$ chains at 140\,000 moves, 5000 moves per $N$:

\begin{center}
\begin{tabular}{rrcccc}
\toprule
$T$ & $N$ & acceptance $a_N$ & $N a_N$ & s/move & s per decorrelation \\
\midrule
32  & 1  & 0.159 & 0.159 & 0.153 & $>384$ \\
32  & 2  & 0.051 & 0.102 & 0.154 & $>385$ \\
32  & 4  & 0.004 & 0.016 & 0.154 & -- \\
32  & $\ge8$ & 0.000 & 0 & 0.156--0.164 & -- \\
\midrule
100 & 1  & 0.497 & 0.497 & 0.153 & 173 \\
100 & 2  & 0.283 & 0.567 & 0.154 & \textbf{149} ($1.16\times$) \\
100 & 4  & 0.135 & 0.541 & 0.154 & 173 \\
100 & 8  & 0.038 & 0.302 & 0.156 & 266 \\
100 & 16 & 0.003 & 0.053 & 0.158 & $>396$ \\
\bottomrule
\end{tabular}
\end{center}

\noindent
The cost of a move barely depends on $N$ (0.153\,s at $N=1$, 0.164\,s at
$N=32$), so the problem is acceptance: $a_N\approx a_1^N$. The reason is that
each site has its own uphill cost (median $\Delta U\approx222$ for one
mutation near the reference, far above $T$), and an $N$-site move must pay all
$N$ costs in one accept/reject. Site--site coupling is minor (13\% of the
summed single-site effect for dispersed pairs, 27\% within 30 residues). With
$a_N=a_1^N$ the best $N$ is $\approx-1/\ln a_1$, so $N=1$ wins whenever
$a_1<0.37$, i.e.\ at every $T$ that matters. Adaptive $N$ or ``choose weakly
coupled sites'' cannot beat the best fixed $N$.

\emph{Answer to the original question:} one mutation per step at and below the
crossover; two at high $T$, for 16\%.

\subsection{Replica exchange (parallel tempering) pilot: no gain at this length}

16 temperatures from 12 to 316, 2000 sweeps (2.45\,s per sweep). Swap
acceptance between neighbours was 0.18--0.47, so the ladder works, and walkers
travelled across it. But each rung was no closer to equilibrium than a plain
chain after the same number of moves: swaps only help once the hot rungs are
equilibrated, which here takes $\sim$20\,000 sweeps (14\,h per ladder, at
$16\times$ the cost per sweep). For three temperatures, plain chains are
cheaper. Code kept (\file{run\_pt.py}) for a temperature scan later.

\subsection{What the faster sampler has produced}

\begin{itemize}
\item \textbf{Crossover location.} $\Hd$ (Hamming distance to the reference)
  falls smoothly with $T$: $\sim$380 at $T=12$, $>330$ at $T=18$, $\sim$495 at
  $T=32$, $\sim$528 at $T=100$ (random sequence: 538). A one-coordinate free
  energy $F(h)=U(h)-T\ln[\binom{L}{h}19^h]$ gives
  $T^*(h)=(dU/dh)/\ln[19(L-h)/h]\approx7$--15 at $h\sim300$, consistent with a
  gradual crossover around $T\approx8$--15 rather than a sharp transition.
\item \textbf{Low $T$.} $T=12$ reference chain: 580\,000 moves, $\Hd$ plateau
  $\sim$380, $U\approx6.9\times10^3$, acceptance 0.05--0.06; being extended
  to $10^6$ (ends today $\sim$10:00). $T=5$ queued next. Random starts never
  reach the reference basin below the crossover (stuck at $\Hd\approx530$),
  so $\hat R$ between reference and random starts cannot be met there.
\end{itemize}

% ===========================================================================
\section{What limits speed now}

\begin{enumerate}
\item \textbf{Acceptance at low $T$.} 0.5--4\% at $T=8$--15 at the start,
  5--6\% at the $T=12$ plateau. More than 94\% of the GPU time at low $T$ is
  spent on rejected moves.
\item \textbf{The gradient is a weak guide at $L=566$.} Its predicted sign of
  $\Delta U$ is right for 13 of 20 single mutations, and its top residue is
  the true best at 2 of 10 sites. The proposal is informed, but not by much.
\item \textbf{One chain does not fill the GPU (estimate).} One forward is
  $\sim$60\,ms for $\sim$1.6\,TFLOP, about 15--20\% of the card's bf16 peak.
  Not yet measured; this decides between the options below.
\end{enumerate}

% ===========================================================================
\section{Options from here}

Gains are wall-clock accepted changes per second relative to today's one-pass
move. $t_{19}$ = time of one batched forward over 19 sequences, not yet
measured (between $\sim$0.25\,s with good batching and $\sim$1.1\,s with none).

\begin{center}
\small
\begin{tabular}{R{3.1cm}R{5.2cm}R{3.6cm}R{2.6cm}}
\toprule
option & idea & expected gain & status / effort \\
\midrule
\textbf{A.} Metropolized Gibbs (batched exact single-site move)
 & at a uniform site, exact $U$ of all 19 alternatives in one batched forward;
   propose $j$ with prob.\ $\pi_j/(1-\pi_{\mathrm{cur}})$, accept
   $\min[1,(1-\pi_{\mathrm{cur}})/(1-\pi_j)]$
 & $T=8$: $\sim12\times$ (good batching) or $2.6\times$ (none);
   $T=32$: $3.8\times$ / $0.8\times$; $T=100$: $1.2\times$ / $0.3\times$
 & $\sim$1 day + toy Boltzmann test; needs $t_{19}$ \\
\addlinespace
\textbf{B.} Parallel chains in one batch
 & $B$ chains share each forward and backward
 & total throughput $B\,t_1/t_B$, not a faster single chain; helps 4-chain
   $\hat R$ runs and PT
 & $\sim$1 day; $B\le3$ on 24\,GB \\
\addlinespace
\textbf{C.} CUDA graphs / \texttt{torch.compile}
 & same move, less overhead
 & unknown, plausibly 1.3--2$\times$; stacks with A and B
 & 0.5--2 days \\
\addlinespace
\textbf{D.} Gradient-chosen site (Gibbs-with-gradients)
 & choose site and residue from all $L\times19$ predicted $\Delta U$ (already
   cached by one-pass)
 & small while the gradient ranks near chance
 & 0.5 day \\
\addlinespace
\textbf{E.} More sites per move
 & adaptive $N$, weakly coupled sites
 & $\le1.16\times$, only at high $T$
 & not recommended \\
\addlinespace
\textbf{F.} Parallel tempering
 & production ladder with fewer rungs
 & useful for a $T$ scan, not for 2--3 temperatures
 & code ready \\
\bottomrule
\end{tabular}
\end{center}

\paragraph{The benchmark that decides.} \file{code/tests/bench\_batch.py}
is written but not yet run ($\sim$10 minutes on astaroth). It measures forward and
forward+backward time for batches of 1--32 sequences (gives $t_{19}$ and the
parallel-chain throughput), the fraction of a forward during which the GPU is
running kernels, and, at the $T=12$ state, the exact $\Delta U$ of all 19
mutants at 20 sites (gives the gradient's ranking quality and the predicted
Metropolized-Gibbs gain for $T=8$--100). It is on the branch \file{claude/project-thread-uu4fbg}. Decision rules:
\begin{itemize}
\item $t_{19}\lesssim0.3$\,s (about $5\times$ one forward): build A first.
\item batching barely helps but the GPU is mostly idle: C first.
\item gradient--exact Spearman $>0.5$: keep the gradient as the site selector
  inside A (a hybrid of A and D).
\end{itemize}

\paragraph{Keeping the gradient.} Option A needs no gradient: the reverse move
uses the same 20 exact energies. The gradient-informed move is the subject of
the thesis, so there are three ways to frame A:
\begin{enumerate}
\item A as a baseline: the exact local move against which the gradient move's
  informedness is measured (how close does the one-pass proposal get to the
  ideal single-site proposal?).
\item A hybrid: the gradient picks \emph{which site} to update (option D),
  A picks the residue exactly.
\item A as the production sampler for equilibrium data at low $T$, with the
  gradient move kept for the method study.
\end{enumerate}

% ===========================================================================
\section{Beyond speed: the structural check}

The plan was to check that low-$T$ samples keep the native fold
(Rg, pLDDT from ESM3 structure decoding) and that low-entropy sites match
sites conserved in nature. Both checks are inconclusive for HA0, for reasons
outside the sampler:
\begin{itemize}
\item \textbf{The decoder does not fold the reference.} ESM3-open decodes the
  reference to pLDDT 0.51/0.60/0.62 and Rg 45/54/70\,\AA\ at 8/32/128 steps
  (0.61 at 200, 0.59 at 256). The ectodomain alone (16--528) is no better
  (pLDDT 0.51--0.63). The decode is deterministic, so the spread over $T=12$
  samples (Rg 42--65\,\AA) is real but not a fold readout. HA is a trimer, and
  ESM3 predicts it as a monomer.
\item \textbf{The H3N2 alignment is too narrow.} 457 of 566 sites are below
  0.1 bits; the variable ones are mostly antigenic. Spearman between sampled
  and H3N2 site entropy: 0.07 at $T=12$, $\approx0$ at $T=32$, 100. Overlap of
  the 10\% most conserved sites: 8 vs 5.7 expected ($p=0.2$). A multi-subtype
  HA alignment is being built.
\end{itemize}
Options for the structure: a multimer predictor (AlphaFold2-multimer via
ColabFold, Boltz-1, Chai-1) run on exported sequences as a trimer, or ESM3 with
the three chains joined by a 25-glycine linker as a quick test. Caveat: reusing
the reference MSA for sampled sequences at $\sim$33\% identity would bias
predictions toward the native fold; single-sequence or MSA-free modes avoid
this but are weaker.

% ===========================================================================
\section{Questions for the supervisors}

\begin{enumerate}
\item Is it acceptable to use an exact batched move (Metropolized Gibbs) with
  no gradient in the proposal, as a baseline, a hybrid, or the production
  sampler at low $T$? (Pending the benchmark numbers.)
\item The $\Hd(T)$ curve is a smooth crossover. Which temperatures should the
  final study use (current: 5, 12, 32, 100)?
\item Below the crossover, random starts never reach the reference basin. Is
  ``a long reference-started chain with a flat $\Hd$ and $U$'' an acceptable
  equilibrium criterion there, in place of $\hat R$ between independent
  starts?
\item Is HA0 the right test case, given that it is a trimer that the
  structure decoder cannot fold? A monomeric reference (e.g.\ protein~G,
  already used) would make the structural check meaningful.
\item Which structure predictor, and is the MSA bias acceptable?
\end{enumerate}

% ===========================================================================
\section*{Immediate next steps}
\begin{itemize}
\item Run \file{bench\_batch.py} (10 min) between the end of $T=12$ and the start of $T=5$.
\item $T=12$ to $10^6$ moves (ends today $\sim$10:00); rerun structure and site
  entropy comparison on the last 500\,000 moves with the multi-subtype HA
  alignment; then resume $T=5$.
\item Build option A (or C) according to the benchmark and the answer to
  question~1.
\end{itemize}

\end{document}
